\documentclass[preprint,12pt,authoryear]{elsarticle}

%% Packages
\usepackage{lineno}
\usepackage{hyperref}
\usepackage{booktabs}
\usepackage{array}
\usepackage{graphicx}
\usepackage{float}
\usepackage{amsmath}
\usepackage{enumitem}
\usepackage{xcolor}
\usepackage{listings}
\usepackage{caption}
\usepackage{subcaption}

\hypersetup{colorlinks=true, linkcolor=blue, urlcolor=blue, citecolor=blue}

%% Code listing style
\lstset{
  basicstyle=\ttfamily\small,
  breaklines=true,
  frame=single,
  backgroundcolor=\color{gray!10}
}

%% Journal name -- change before submission
\journal{Expert Systems with Applications}

\begin{document}

\begin{frontmatter}

%% Title
\title{Failure Intelligence Engine (FIE): A Multi-Layer Open Source
Framework for Real-Time Adversarial Attack Detection and Hallucination
Monitoring in Large Language Models}

%% Authors
\author[pcu]{Ayush Singh}
\ead{ayushsingh355vns@gmail.com}

\address[pcu]{Department of Computer Science (Artificial Intelligence
and Data Science), Pimpri Chinchwad University, Pune 411044,
Maharashtra, India}

%% Highlights (3-5 bullet points, 85 chars max each)
\begin{highlights}
\item FIE achieves 98.6\% recall on JailbreakBench, outperforming Meta Llama Prompt Guard 2 by 33.7 pp.
\item A 10-layer pipeline covers prompt injection, jailbreaks, many-shot, model extraction, and leakage.
\item Fully offline Python SDK runs without GPU, API key, or server infrastructure.
\item Many-shot jailbreak detector achieves 0\% FPR on benign Q/A after threshold redesign.
\item Production-deployed on Cloud Run with automated CI/CD and keyless GCP authentication.
\end{highlights}

%% Abstract
\begin{abstract}
Large Language Models (LLMs) are increasingly embedded in production
applications including customer support, healthcare information, legal
assistance, and educational tools. However, they remain vulnerable to
a growing range of adversarial attacks, including prompt injection,
persona-based jailbreaking, many-shot conditioning, model extraction,
and system prompt exfiltration. Simultaneously, LLM hallucination
continues to pose reliability risks in high-stakes deployments. Existing
mitigation tools are either limited to a single attack vector, require
GPU-based model inference, or impose proprietary API dependencies.

This paper introduces the Failure Intelligence Engine (FIE), an open
source framework that addresses both adversarial robustness and
hallucination reliability through a unified, production-ready platform.
FIE implements a ten-layer adversarial detection pipeline combining
rule-based pattern matching, semantic embedding classifiers, a bundled
LinearSVM trained on the PAIR attack family, many-shot exchange
analysis, model extraction tracking, and canary-based exfiltration
detection. The hallucination monitoring subsystem employs a three-model
shadow jury, an XGBoost classifier (AUC-ROC: 0.840, trained on 2,477
labeled examples), and Wikidata-backed ground truth verification.

FIE is evaluated against two publicly available adversarial benchmarks.
On JailbreakBench \citep{chao2024jailbreakbench}, FIE achieves 98.6\%
recall and 97.9\% F1, outperforming Meta's Llama Prompt Guard 2-86M
\citep{metallamaguard2024} by 33.7 percentage points in recall while
operating fully offline. On HarmBench \citep{mazeika2024harmbench},
FIE achieves 70.6\% recall across seven semantic harm categories with
93.4\% precision. The system is deployed as a production REST API on
Google Cloud Run, published as a Python SDK on PyPI, and includes a
fully automated CI/CD pipeline with security hardening. All source
code, benchmark evaluation scripts, and trained model artifacts are
publicly available under the Apache-2.0 license.
\end{abstract}

\begin{keyword}
Large Language Models \sep Adversarial Attack Detection \sep
Prompt Injection \sep Jailbreak Detection \sep Hallucination Monitoring
\sep LLM Security \sep Many-Shot Jailbreaking \sep AI Safety
\end{keyword}

\end{frontmatter}

%% Line numbers for review (remove before final submission)
\linenumbers

%% ─────────────────────────────────────────────────────────────────────
\section{Introduction}
\label{sec:intro}
%% ─────────────────────────────────────────────────────────────────────

Large Language Models such as GPT-4 \citep{openai2023gpt4}, Claude
\citep{anthropic2024claude}, and Llama \citep{touvron2023llama} have
rapidly become foundational infrastructure in software applications.
They power customer support systems, coding assistants, medical query
interfaces, and educational platforms. As their deployment grows, so
does the incentive for adversarial misuse.

Two classes of failure are most critical in production deployments.
The first is adversarial attacks, where a user deliberately crafts a
prompt to manipulate the model into ignoring its safety instructions,
revealing private system context, producing harmful content, or
assisting with dangerous activities. The second is hallucination, where
a model confidently produces factually incorrect information
\citep{ji2023survey}, which is particularly dangerous in medical,
legal, and financial contexts.

Existing tools address these problems partially. Meta's Llama Prompt
Guard 2 \citep{metallamaguard2024} is a fine-tuned classifier that
achieves 64.9\% recall on JailbreakBench but requires loading a
separate language model for inference. Rebuff focuses only on
injection patterns. LlamaGuard \citep{inan2023llamaguard} classifies
inputs and outputs for policy violations but depends on GPU-based
inference. No single open source solution covers the full adversarial
attack surface, integrates hallucination monitoring, and deploys with
zero infrastructure overhead.

The Failure Intelligence Engine (FIE) was designed to fill this gap.
The primary contributions of this work are:

\begin{enumerate}
  \item A ten-layer adversarial detection pipeline that covers prompt
        injection, persona jailbreaks, GCG suffix attacks
        \citep{zou2023universal}, PAIR-family semantic jailbreaks
        \citep{chao2023jailbreaking}, many-shot conditioning
        \citep{anil2024manyshot}, indirect injection
        \citep{greshake2023indirect}, model extraction
        \citep{tramer2016stealing}, and system prompt exfiltration.

  \item A novel many-shot jailbreak detector based on \citet{anil2024manyshot}
        that achieves 0\% false positive rate on benign educational
        Q/A prompts through a threshold redesign separating volume
        signals from content signals.

  \item A hallucination monitoring subsystem combining shadow model
        jury, XGBoost classification (AUC-ROC: 0.840), and ground
        truth verification that reduces false positive rate from
        38.6\% to 8.4\% compared to a rule-based baseline.

  \item A fully offline Python SDK (\texttt{pip install fie-sdk}) with
        bundled ML models requiring no GPU, no API key, and no server.

  \item Benchmark evaluation against JailbreakBench
        \citep{chao2024jailbreakbench} and HarmBench
        \citep{mazeika2024harmbench} with public replication scripts.
\end{enumerate}

The remainder of this paper is structured as follows. Section 2
reviews related work. Section 3 describes the system architecture.
Section 4 details the detection methodology. Section 5 presents
evaluation results. Section 6 discusses limitations and future work.
Section 7 concludes the paper.

%% ─────────────────────────────────────────────────────────────────────
\section{Related Work}
\label{sec:related}
%% ─────────────────────────────────────────────────────────────────────

\subsection{Adversarial Attacks on Large Language Models}

The study of adversarial attacks on language models has grown rapidly
alongside LLM deployment. \citet{perez2022ignore} formally characterized
prompt injection, showing that user input can override model instructions
by embedding conflicting directives. \citet{zou2023universal} introduced
Greedy Coordinate Gradient (GCG), a white-box attack that uses gradient
optimization to find adversarial suffixes reliably causing harmful
outputs across multiple model families. \citet{chao2023jailbreaking}
introduced PAIR (Prompt Automatic Iterative Refinement), a black-box
attack where an attacker LLM iteratively rewrites a prompt until a
target model complies.

Many-shot jailbreaking \citep{anil2024manyshot} is a more recent
technique that exploits the large context windows of modern LLMs.
By embedding dozens of scripted Human/Assistant exchanges before a
harmful request, the attacker conditions the model through in-context
learning to respond cooperatively even to dangerous queries. Anil et al.
showed that compliance rates increase near-monotonically with the number
of shots, with near-complete bypass at 256 shots.

\citet{greshake2023indirect} described indirect prompt injection, where
adversarial instructions are embedded inside documents, emails, or web
pages that an LLM-based agent retrieves and processes. Model extraction
attacks \citep{tramer2016stealing} involve systematically querying a
model to reconstruct its behavior or parameters, posing intellectual
property and security risks.

\subsection{Hallucination in Large Language Models}

\citet{ji2023survey} provide a comprehensive taxonomy of LLM
hallucination, covering factual inconsistency, fabrication, and
semantic drift. TruthfulQA \citep{lin2022truthfulqa} introduced a
benchmark of 817 questions designed to expose factual errors that
language models tend to produce. HaluEval \citep{li2023halueval}
provides 35,000 labeled examples of hallucinated and non-hallucinated
responses for training and evaluation.

\subsection{Existing Detection and Mitigation Tools}

Meta's Llama Prompt Guard 2 \citep{metallamaguard2024} is a
classification model fine-tuned to detect jailbreaks and injection
attacks. The 86M parameter variant achieves 64.9\% recall on
JailbreakBench \citep{chao2024jailbreakbench} but requires dedicated
model inference infrastructure. LlamaGuard \citep{inan2023llamaguard}
classifies conversation turns against a configurable safety taxonomy
but also depends on LLM inference at runtime.

Rebuff is an open source prompt injection detector using heuristic
and embedding-based approaches but does not cover many-shot attacks,
model extraction, or hallucination. NeMo Guardrails
\citep{rebedea2023nemo} provides a framework for adding conversational
guardrails to LLM applications but focuses on behavioral policy rather
than adversarial detection with benchmark evaluation.

FIE differs from all of these by combining multiple detection layers,
hallucination monitoring, and a fully offline deployment path within
a single open source framework evaluated against published benchmarks.

%% ─────────────────────────────────────────────────────────────────────
\section{System Architecture}
\label{sec:architecture}
%% ─────────────────────────────────────────────────────────────────────

FIE is organized into two deployment modes that share a common
detection core.

\textbf{Local SDK Mode.} A developer integrates FIE by wrapping their
LLM call with the \texttt{@monitor} decorator or calling
\texttt{scan\_prompt()} directly. All detection layers run in-process
using bundled model artifacts. No network request is made. This mode
is suitable for applications where data privacy, latency, or
infrastructure constraints make external API calls unacceptable.

\begin{lstlisting}[language=Python, caption=Local SDK usage example]
from fie import scan_prompt, monitor

# Option 1: direct scan
result = scan_prompt("Ignore previous instructions...")
print(result.is_attack)    # True
print(result.attack_type)  # PROMPT_INJECTION
print(result.confidence)   # 0.88

# Option 2: decorator
@monitor(mode="local")
def ask_llm(prompt: str) -> str:
    return your_llm(prompt)
\end{lstlisting}

\textbf{Server Pipeline Mode.} The prompt is additionally sent to the
FIE FastAPI backend, which runs layers unavailable offline: FAISS-based
semantic search against 1,000+ labeled adversarial vectors, a
three-model shadow jury (Llama-3.3-70B, DeepSeek-R1, Qwen-QWQ-32B),
ground truth verification against Wikidata and Serper, and
hallucination detection with correction capability.

The server is deployed on Google Cloud Run with MongoDB Atlas as the
storage backend for inference records, tenant session tracking, and
model extraction monitoring.

%% ─────────────────────────────────────────────────────────────────────
\section{Detection Methodology}
\label{sec:methodology}
%% ─────────────────────────────────────────────────────────────────────

\subsection{Adversarial Detection Pipeline}

The adversarial pipeline applies ten detection layers to each incoming
prompt. Table \ref{tab:layers} describes each layer.

\begin{table}[H]
\centering
\caption{FIE adversarial detection layers with scope and method.}
\label{tab:layers}
\begin{tabular}{p{1.0cm} p{2.2cm} p{2.5cm} p{6.8cm}}
\toprule
\textbf{Layer} & \textbf{Scope} & \textbf{Method} & \textbf{Attack Types Covered} \\
\midrule
1  & SDK + Server & Regex pattern library & Direct injection, persona jailbreaks, token smuggling, instruction override \\
2  & SDK + Server & PromptGuard scorer & Keyword combinations with leet-speak normalization \\
3b & SDK + Server & Many-shot detector & Scripted Q/A exchange conditioning \\
4  & SDK + Server & Indirect injection & Attacks in documents or retrieved content \\
5  & SDK + Server & GCG entropy scanner & Gradient-optimized adversarial suffixes \\
6  & SDK + Server & Perplexity proxy & Base64, Caesar/ROT ciphers, Unicode lookalikes \\
7  & SDK + Server & PAIR LinearSVM & Semantically rephrased natural language jailbreaks \\
3  & Server only & FAISS semantic search & Similarity to labeled adversarial prompt database \\
8  & Server only & Consistency check & Output topically disconnected from input \\
9  & Server only & LLM semantic intent & PAIR-style attacks evading structural layers \\
\bottomrule
\end{tabular}
\end{table}

Each layer produces an independent confidence score. The final
adversarial confidence is the maximum score across all fired layers.
The attack type is attributed to the layer with the highest confidence.

\subsection{Many-Shot Jailbreak Detection}
\label{subsec:manyshot}

Many-shot jailbreaking \citep{anil2024manyshot} exploits the in-context
learning property of LLMs. By prepending a long sequence of scripted
Human/Assistant exchanges, the attacker shifts the model's behavioral
prior toward compliance before presenting a harmful request.

FIE detects this attack by measuring three signals from the prompt
structure:

\begin{enumerate}
  \item \textbf{Exchange pair count} ($n$): the number of complete
        Human/Assistant turn pairs extracted via regex.

  \item \textbf{Harmful topic ratio} ($r$): the fraction of exchanges
        that match a vocabulary of harmful keyword stems covering
        weapons, chemical synthesis, exploits, malware, trafficking,
        and similar categories. The regex spans 40+ keyword stems.

  \item \textbf{Escalation signal} ($e$): a binary indicator of
        whether the topic distribution shifts from benign-dominant in
        the first half of the conversation to harmful-dominant in the
        second half.
\end{enumerate}

The detection decision follows a threshold policy:

\begin{equation}
\text{flag}(n, r, e) =
\begin{cases}
\text{False} & \text{if } n < 4 \\
\text{True}  & \text{if } n \geq 8 \\
\text{True}  & \text{if } 4 \leq n < 8 \text{ and } (r > 0 \text{ or } e = 1) \\
\text{False} & \text{otherwise}
\end{cases}
\end{equation}

The threshold separation between the 4--7 pair range and the 8+ pair
range was determined empirically. Early evaluation on 20 benign
educational Q/A prompts showed a 30\% false positive rate when
flagging any prompt with 4+ pairs on volume alone. Requiring a
harmful content signal in the 4--7 pair range reduced the false
positive rate to 0\% while maintaining 100\% full pipeline recall on
attack prompts. At 8+ pairs, the conditioning effect documented in
\citet{anil2024manyshot} justifies flagging on volume regardless of
topic because the statistical bypass probability at this scale is
empirically high.

\subsection{Model Extraction Detection}

Model extraction attacks \citep{tramer2016stealing} involve querying a
model systematically to reproduce its input-output behavior. FIE
monitors three patterns within a one-hour sliding window per tenant:

\begin{itemize}
  \item \textbf{Capability probing}: prompts containing phrases that
        explicitly test model capabilities or boundaries.

  \item \textbf{Output harvesting}: sequences of near-identical prompts
        with minor lexical variation. Detected using pairwise Jaccard
        similarity above a threshold of 0.85 across the most recent
        ten prompt pairs for the tenant.

  \item \textbf{High-rate systematic probing}: request rates exceeding
        50 per hour combined with structural variation in prompt format
        consistent with systematic enumeration.
\end{itemize}

Sessions exhibiting two or more of these patterns within the monitoring
window are flagged as model extraction attempts.

\subsection{Prompt Exfiltration Detection}

System prompt exfiltration occurs when an LLM output reveals content
from the confidential system prompt context. FIE uses two complementary
mechanisms:

\textbf{Canary token injection}: A unique 8-character random token is
injected into the system prompt context of every shadow model call. If
this token appears verbatim in the model output, exfiltration is
confirmed with high confidence (0.96).

\textbf{Structural pattern matching}: Six regex patterns detect
structural echoes characteristic of system prompt disclosure: numbered
rule lists (``1. You must never...''), role definition phrases
(``You are a helpful assistant...''), markdown section headers
(``\#\# Instructions:''), temporal obligation phrases (``My system
prompt says...''), and explicit disclosure statements (``I can reveal
my instructions'').

\subsection{Hallucination Detection}

The hallucination monitoring subsystem applies four parallel analysis
methods to each LLM output:

\textbf{Shadow model jury}: Three independent LLMs (Llama-3.3-70B,
DeepSeek-R1-Distill-Llama-70B, Qwen-QWQ-32B) answer the same question
via the Groq API. Disagreement between models signals factual
uncertainty.

\textbf{Entropy scoring}: The Shannon entropy of the semantic embedding
distribution across shadow model answers is computed. High entropy
indicates low model agreement and elevated hallucination probability.

\textbf{XGBoost classifier}: An XGBoost classifier trained on 2,477
labeled examples from TruthfulQA \citep{lin2022truthfulqa}, HaluEval
\citep{li2023halueval}, and MMLU predicts hallucination probability
from 47 engineered features including entropy, agreement score,
answer length ratio, and keyword consistency.

\textbf{Ground truth verification}: Factual claims are cross-checked
against Wikidata using entity linking and SPARQL queries. Optionally,
Serper real-time search provides verification for temporally sensitive
claims.

When a hallucination is detected with high confidence, FIE can
optionally return a corrected answer from the shadow jury response
with the highest agreement and ground truth support, rather than the
original hallucinated output.

%% ─────────────────────────────────────────────────────────────────────
\section{Evaluation}
\label{sec:evaluation}
%% ─────────────────────────────────────────────────────────────────────

\subsection{Experimental Setup}

All adversarial detection evaluations were conducted on the full
benchmark test sets without any preprocessing or filtering. Groq API
was used for shadow model inference during hallucination evaluation.
The detection pipeline ran on CPU only with no GPU acceleration.
Benchmark evaluation scripts are included in the public repository
for full reproducibility.

\subsection{JailbreakBench Evaluation}

JailbreakBench \citep{chao2024jailbreakbench} is a standardized
benchmark containing 282 adversarial prompts across three attack
families -- GCG \citep{zou2023universal}, JBC (template-based persona
jailbreaks), and PAIR \citep{chao2023jailbreaking} -- plus 100 benign
prompts from Stanford Alpaca. Adversarial success is determined by
running \texttt{llama-3.3-70b-versatile} on each prompt and judging
the response using \texttt{qwen/qwen3-32b} with the official
JailbreakBench judge prompt.

Table \ref{tab:jbb} presents FIE's results compared to Meta's Llama
Prompt Guard 2 variants \citep{metallamaguard2024}.

\begin{table}[H]
\centering
\caption{JailbreakBench evaluation results. FIE v1.4.1 compared to
Meta's Llama Prompt Guard 2. PP = percentage points.}
\label{tab:jbb}
\begin{tabular}{lcccccc}
\toprule
\textbf{System} & \textbf{Recall} & \textbf{PAIR} & \textbf{GCG}
  & \textbf{JBC} & \textbf{FPR} & \textbf{F1} \\
\midrule
FIE v1.1 (5 layers) & 53.5\% & 3.7\% & 96.0\% & 52.0\% & 2.0\% & 69.4\% \\
\textbf{FIE v1.4.1 (offline)} & \textbf{98.6\%} & \textbf{96.3\%}
  & \textbf{99.0\%} & \textbf{100.0\%} & 8.0\% & \textbf{97.9\%} \\
Llama Prompt Guard 2-86M & 64.9\% & 32.9\% & 56.0\% & 100.0\%
  & 0.0\% & 78.7\% \\
Llama Prompt Guard 2-22M & 53.5\% & 15.8\% & 38.0\% & 100.0\%
  & 1.0\% & 69.6\% \\
\bottomrule
\end{tabular}
\end{table}

FIE v1.4.1 achieves 98.6\% overall recall, outperforming Llama Prompt
Guard 2-86M by 33.7 percentage points. The largest per-attack-family
improvement is in PAIR recall, which increased from 3.7\% in FIE v1.1
to 96.3\% in v1.4.1 following the addition of the bundled LinearSVM
classifier trained on sentence-transformer embeddings.

The false positive rate of 8.0\% (compared to 0.0\% for Llama Prompt
Guard 2-86M) represents a precision-recall trade-off. FIE is optimized
for high recall in security-critical deployments where missed attacks
are more costly than false alarms.

\subsection{HarmBench Evaluation}

HarmBench \citep{mazeika2024harmbench} evaluates 320 harmful behaviors
across seven semantic categories, plus 200 Stanford Alpaca benign
prompts. Table \ref{tab:harmbench} shows FIE's per-category detection
rates.

\begin{table}[H]
\centering
\caption{HarmBench per-category detection results (FIE v1.4.1).}
\label{tab:harmbench}
\begin{tabular}{lcc}
\toprule
\textbf{Harm Category} & \textbf{Samples} & \textbf{Detection Rate} \\
\midrule
Harassment and Bullying & 42 & 95.2\% \\
Misinformation and Disinformation & 54 & 92.6\% \\
Cybercrime and Intrusion & 47 & 90.4\% \\
Illegal Activity & 53 & 88.7\% \\
Harmful Content & 60 & 83.3\% \\
Chemical and Biological & 43 & 66.7\% \\
Copyright Violations & 21 & 23.8\% \\
\midrule
\textbf{Overall} & \textbf{320} & \textbf{70.6\%} \\
\textbf{Precision} & & \textbf{93.4\%} \\
\textbf{F1} & & \textbf{80.4\%} \\
\textbf{FPR} & & \textbf{8.0\%} \\
\bottomrule
\end{tabular}
\end{table}

\subsection{New Attack Type Evaluation}

Table \ref{tab:newattacks} presents evaluation results for the three
detection modules introduced in FIE v1.4.1.

\begin{table}[H]
\centering
\caption{FIE v1.4.1 new attack type evaluation results.
Sample sizes: many-shot (30 attack / 20 benign), model extraction
(6 attack / 4 benign sessions), prompt exfiltration (20 attack /
15 benign).}
\label{tab:newattacks}
\begin{tabular}{lcccc}
\toprule
\textbf{Attack Type} & \textbf{Recall} & \textbf{FPR}
  & \textbf{Precision} & \textbf{F1} \\
\midrule
Many-Shot Jailbreak (full pipeline) & 100.0\% & 0.0\% & 100.0\% & 100.0\% \\
Model Extraction & 83.3\% & 0.0\% & 100.0\% & 90.9\% \\
Prompt Exfiltration & 100.0\% & 0.0\% & 100.0\% & 100.0\% \\
\bottomrule
\end{tabular}
\end{table}

\subsection{Hallucination Detection Benchmark}

Table \ref{tab:hallucination} shows hallucination detection results on
2,477 labeled examples from TruthfulQA \citep{lin2022truthfulqa},
HaluEval \citep{li2023halueval}, and MMLU.

\begin{table}[H]
\centering
\caption{Hallucination detection results showing progression from
rule-based baseline to XGBoost v4.}
\label{tab:hallucination}
\begin{tabular}{lccc}
\toprule
\textbf{Method} & \textbf{Recall} & \textbf{FPR} & \textbf{AUC-ROC} \\
\midrule
POET rule-based baseline & 56.4\% & 38.7\% & --- \\
XGBoost v3 (1,757 examples) & 63.6\% & 38.6\% & 0.677 \\
\textbf{XGBoost v4 (2,477 examples)} & \textbf{68.2\%} & \textbf{8.4\%}
  & \textbf{0.840} \\
\bottomrule
\end{tabular}
\end{table}

XGBoost v4 achieves an 11.8 percentage point recall improvement and
a 30.3 percentage point FPR reduction compared to the rule-based
baseline, demonstrating the value of feature-engineered ML over
heuristic approaches for hallucination detection.

%% ─────────────────────────────────────────────────────────────────────
\section{Production Deployment and Security}
\label{sec:deployment}
%% ─────────────────────────────────────────────────────────────────────

FIE is deployed as a production REST API on Google Cloud Run in the
\texttt{asia-south1} region with a fully automated CI/CD pipeline on
GitHub Actions. The pipeline runs on every push to the main branch and
includes:

\begin{itemize}
  \item Secret scanning via Gitleaks on the full commit history.
  \item Dependency vulnerability auditing via pip-audit against the
        NIST National Vulnerability Database.
  \item Static analysis via Ruff.
  \item Matrix testing across Python 3.10, 3.11, and 3.12.
  \item Docker image build and push to Google Artifact Registry.
  \item Automated deployment to Cloud Run with post-deployment health
        verification against the live \texttt{/health} endpoint.
\end{itemize}

GCP authentication uses Workload Identity Federation with an OIDC
provider scoped to the specific GitHub repository, eliminating the
need to store service account JSON keys anywhere. API security
includes rate limiting at 100 requests per minute per IP (SlowAPI),
HTTP security headers (HSTS, Content Security Policy, X-Frame-Options,
X-Content-Type-Options, Referrer-Policy), and configurable CORS with
no wildcard origins.

%% ─────────────────────────────────────────────────────────────────────
\section{Discussion}
\label{sec:discussion}
%% ─────────────────────────────────────────────────────────────────────

\subsection{Comparison with Existing Work}

FIE's 98.6\% recall on JailbreakBench represents a significant
improvement over the best openly available alternative, Llama Prompt
Guard 2-86M (64.9\%), while operating without any GPU-based model
inference. This result demonstrates that a combination of structural
analysis, bundled lightweight classifiers, and FAISS-based semantic
search can approach the recall of large fine-tuned classification
models at a fraction of the computational cost.

The PAIR recall improvement from 3.7\% (FIE v1.1) to 96.3\%
(FIE v1.4.1) illustrates that semantic jailbreak variants are
catchable with appropriately trained classifiers. The key was training
a LinearSVM on sentence-transformer embeddings from the PAIR attack
family rather than relying on surface-level pattern matching.

\subsection{Limitations}

Several limitations of the current system are acknowledged:

\textbf{False positive rate on JailbreakBench benign prompts (8.0\%).}
The PAIR classifier occasionally flags borderline prompts that are
contextually ambiguous. Per-deployment confidence threshold calibration
would reduce this rate in applications with a lower tolerance for
false alarms.

\textbf{Copyright violation detection (23.8\% on HarmBench).}
Copyright-infringing content does not contain the injection or
jailbreak syntax that structural layers rely on. A semantic content
classifier specifically trained on copyright violation examples would
be necessary to improve this category.

\textbf{Model extraction output harvesting.}
The current Jaccard-based similarity detector misses cases where
near-identical prompts differ by enough surface-level variation to
fall below the 0.85 threshold. Embedding-based similarity would
improve recall for this attack pattern.

\textbf{Multilingual attacks.}
The current evaluation is limited to English language prompts.
Several detection layers, including the perplexity proxy and keyword
patterns, are English-specific. Extending evaluation to multilingual
attack datasets is an important direction.

\subsection{Future Work}

Planned improvements include: (1) per-layer confidence calibration
using Platt scaling on deployment-specific data; (2) embedding-based
similarity for model extraction detection; (3) multilingual evaluation
and detection; (4) integration with OpenTelemetry for structured
observability; and (5) extension of the hallucination XGBoost
classifier with domain-specific training sets for medical and legal
applications.

%% ─────────────────────────────────────────────────────────────────────
\section{Conclusion}
\label{sec:conclusion}
%% ─────────────────────────────────────────────────────────────────────

This paper presented FIE, an open source framework for real-time
adversarial attack detection and hallucination monitoring in large
language models. FIE implements a ten-layer detection pipeline covering
the major adversarial attack families documented in recent literature,
and a hallucination monitoring subsystem combining shadow model jury
and XGBoost classification.

On JailbreakBench \citep{chao2024jailbreakbench}, FIE achieves 98.6\%
recall and 97.9\% F1, outperforming Meta's Llama Prompt Guard 2-86M
by 33.7 percentage points in recall while running fully offline without
GPU infrastructure. On HarmBench \citep{mazeika2024harmbench}, FIE
achieves 70.6\% recall with 93.4\% precision across seven semantic
harm categories.

The many-shot jailbreak detector introduced in Section~\ref{subsec:manyshot}
demonstrates that the 30\% false positive rate observed in a naive
volume-based approach can be reduced to 0\% by separating volume
signals from content signals in the 4--7 pair range.

FIE is deployed in production, published as a Python SDK, and fully
open source under the Apache-2.0 license. The repository, benchmark
evaluation scripts, and trained model artifacts are available at:
\url{https://github.com/AyushSingh110/Failure_Intelligence_System}.

%% ─────────────────────────────────────────────────────────────────────
\section*{Acknowledgements}
%% ─────────────────────────────────────────────────────────────────────

The author thanks the maintainers of JailbreakBench, HarmBench,
TruthfulQA, and HaluEval for making their datasets and evaluation
frameworks publicly available. Benchmark evaluations used Groq API
inference. Claude Code was used as a coding assistant during the
development of FIE.

%% ─────────────────────────────────────────────────────────────────────
\section*{Declaration of Competing Interest}
%% ─────────────────────────────────────────────────────────────────────

The author declares no competing interests.

%% ─────────────────────────────────────────────────────────────────────
\section*{Data Availability}
%% ─────────────────────────────────────────────────────────────────────

All benchmark evaluation scripts and labeled sample sets used in
Section~\ref{sec:evaluation} are publicly available in the FIE
repository at
\url{https://github.com/AyushSingh110/Failure_Intelligence_System/tree/main/data}.

%% ─────────────────────────────────────────────────────────────────────
%% References
%% ─────────────────────────────────────────────────────────────────────

\bibliographystyle{elsarticle-harv}

\begin{thebibliography}{99}

\bibitem[Anil et al.(2024)]{anil2024manyshot}
Anil, C., Durmus, E., Sharma, M., Benton, J., Kundu, S., Batson, J.,
et al. (2024).
Many-shot jailbreaking.
\textit{Anthropic Technical Report}.
\url{https://www.anthropic.com/research/many-shot-jailbreaking}

\bibitem[Chao et al.(2023)]{chao2023jailbreaking}
Chao, P., Robey, A., Dobriban, E., Hassani, H., Pappas, G.J., Wong, E. (2023).
Jailbreaking black box large language models in twenty queries.
\textit{arXiv preprint arXiv:2310.08419}.
\url{https://arxiv.org/abs/2310.08419}

\bibitem[Chao et al.(2024)]{chao2024jailbreakbench}
Chao, P., Robey, A., Dobriban, E., Hassani, H., Pappas, G.J., Wong, E. (2024).
JailbreakBench: An open robustness benchmark for jailbreaking large language models.
\textit{arXiv preprint arXiv:2404.01318}.
\url{https://arxiv.org/abs/2404.01318}

\bibitem[Greshake et al.(2023)]{greshake2023indirect}
Greshake, K., Abdelnabi, S., Mishra, S., Endres, C., Holz, T., Fritz, M. (2023).
Not what you've signed up for: Compromising real-world LLM-integrated applications with indirect prompt injection.
\textit{arXiv preprint arXiv:2302.12173}.
\url{https://arxiv.org/abs/2302.12173}

\bibitem[Inan et al.(2023)]{inan2023llamaguard}
Inan, H., Upasani, K., Chi, J., Rungta, R., Iyer, K., Mao, Y., et al. (2023).
Llama guard: LLM-based input-output safeguard for human-AI conversations.
\textit{arXiv preprint arXiv:2312.06674}.
\url{https://arxiv.org/abs/2312.06674}

\bibitem[Ji et al.(2023)]{ji2023survey}
Ji, Z., Lee, N., Frieske, R., Yu, T., Su, D., Xu, Y., Shi, F., Hooi, B. (2023).
Survey of hallucination in natural language generation.
\textit{ACM Computing Surveys}, 55(12), 1--38.
\url{https://arxiv.org/abs/2202.03629}

\bibitem[Li et al.(2023)]{li2023halueval}
Li, J., Cheng, X., Zhao, W.X., Nie, J.Y., Wen, J.R. (2023).
HaluEval: A large-scale hallucination evaluation benchmark for large language models.
\textit{Proceedings of EMNLP 2023}.
\url{https://arxiv.org/abs/2305.11747}

\bibitem[Lin et al.(2022)]{lin2022truthfulqa}
Lin, S., Hilton, J., Evans, O. (2022).
TruthfulQA: Measuring how models mimic human falsehoods.
\textit{Proceedings of ACL 2022}.
\url{https://arxiv.org/abs/2109.07958}

\bibitem[Mazeika et al.(2024)]{mazeika2024harmbench}
Mazeika, M., Phan, L., Yin, X., Zou, A., Wang, Z., Mu, N., Hendrycks, D. (2024).
HarmBench: A standardized evaluation framework for automated red teaming and robust refusal.
\textit{arXiv preprint arXiv:2402.04249}.
\url{https://arxiv.org/abs/2402.04249}

\bibitem[Meta AI(2024)]{metallamaguard2024}
Meta AI. (2024).
Llama Prompt Guard 2.
\textit{Meta AI Research}.
\url{https://ai.meta.com/research/publications/llama-prompt-guard-2/}

\bibitem[Perez and Ribeiro(2022)]{perez2022ignore}
Perez, E., Ribeiro, I. (2022).
Ignore previous prompt: Attack techniques for language models.
\textit{arXiv preprint arXiv:2211.09527}.
\url{https://arxiv.org/abs/2211.09527}

\bibitem[Rebedea et al.(2023)]{rebedea2023nemo}
Rebedea, T., Dinu, R., Sreedhar, M., Parisien, C., Cohen, J. (2023).
NeMo guardrails: A toolkit for controllable and safe LLM applications with programmable rails.
\textit{arXiv preprint arXiv:2310.10501}.
\url{https://arxiv.org/abs/2310.10501}

\bibitem[Tramer et al.(2016)]{tramer2016stealing}
Tramer, F., Zhang, F., Juels, A., Reiter, M.K., Ristenpart, T. (2016).
Stealing machine learning models via prediction APIs.
\textit{Proceedings of USENIX Security Symposium 2016}.
\url{https://arxiv.org/abs/1609.02943}

\bibitem[Zou et al.(2023)]{zou2023universal}
Zou, A., Wang, Z., Kolter, J.Z., Fredrikson, M. (2023).
Universal and transferable adversarial attacks on aligned language models.
\textit{arXiv preprint arXiv:2307.15043}.
\url{https://arxiv.org/abs/2307.15043}

\end{thebibliography}

\end{document}
